\par\Needspace{4\baselineskip}
\begin{center}
\SourceLeaf{027b}\pdfbookmark[2]{9.}{article-009}9.
\end{center}
Verlängert man das Lot, das von einem beliebigen Punkt der Ellipse auf die
Apsidenlinie gefällt wird, rückwärts, bis es den Kreis um den Mittelpunkt der Ellipse
mit dem Radius $a$ schneidet, so ist die Neigung des zum Schnittpunkt gehörenden Radius
gegen die Apsidenlinie (in demselben Sinne verstanden wie oben bei der wahren Anomalie)
gleich der exzentrischen Anomalie, wie sich ohne Weiteres aus Gl. IX des vorigen Artikels
ergibt. Ferner ist klar, dass $r\sin v$ der Abstand eines beliebigen Punktes der
Ellipse von der Apsidenlinie ist. Da dieser nach Gl. VIII $= a\cos\varphi\sin E$ ist,
wird er für $E = 90^0$, also im Mittelpunkt der Ellipse, am größten. Dieser größte
Abstand, der $= a\cos\varphi = \frac{p}{\cos\varphi} = \BareRoot{ap}$ ist, heißt
\emph{kleine Halbachse}. Im Brennpunkt der Ellipse, also für $v = 90^0$, ist dieser
Abstand offenbar $= p$, also gleich dem Halbparameter.
